\section{BPS bracket constants: second attack and repair}
\label{sec:bps-bracket-constants-attack-heal}

\providecommand{\Z}{\mathbb Z}
\providecommand{\C}{\mathbb C}
\providecommand{\R}{\mathbb R}
\providecommand{\La}{\Lambda}
\providecommand{\autcor}{\mathfrak g_{\Delta_5}}
\providecommand{\sdim}{\operatorname{sdim}}
\providecommand{\smult}{\operatorname{smult}}
\providecommand{\Prim}{\operatorname{Prim}}

\subsection{Claim attacked}

The attacked claim is stronger than the denominator theorem:
the Borcherds bracket on
$$
\mathcal A_{\mathrm{den}}=\autcor
$$
is claimed to be recoverable as an honest compact
Hall/CoHA/Drinfeld bracket on protected BPS primitives from the
determinant and active-support data alone.  This is false.  The
denominator presentation gives many low-degree algebra constants.  It
does not give compact Hall extension constants.

\subsection{Low-degree active roots}

Write a positive denominator degree as
$$
\beta=c_1\delta_1+c_2\delta_2+c_3\delta_3,\qquad
\operatorname{ht}(\beta)=c_1+c_2+c_3.
$$
The charge corresponding to \((c_1,c_2,c_3)\) is
$$
\gamma(\beta)=(n,l,m)=(c_1,\ c_1+c_2-c_3,\ c_2).
$$
Indeed
$$
\alpha(n,l,m)=n\delta_1+m\delta_2+(n+m-l)\delta_3.
$$
For heights at most three the active denominator degrees are:
$$
\begin{array}{c|c|c|c}
\beta & \gamma(\beta) & (\beta,\beta) & 4nm-l^2\\
\hline
\delta_i & \gamma_i & 2 & -1\\
\delta_i+\delta_j\ (i<j) & (1,2,1),(1,0,0),(0,0,1) & 0 & 0\\
2\delta_i+\delta_j\ (i\neq j) & \hbox{six real charges} & 2 & -1\\
\delta_1+\delta_2+\delta_3 & (1,1,1) & -6 & 3 .
\end{array}
$$
The six charges in the third row are
$$
(2,3,1),\ (1,3,2),\ (2,1,0),\ (1,-1,0),\
(0,1,2),\ (0,-1,1),
$$
all with discriminant \(-1\).  The Jacobi coefficients used here are
$$
f(0,\pm1)=1,\qquad f(0,0)=10,\qquad
f(1,\pm2)=10,\qquad f(1,\pm1)=-64.
$$
Thus the signed denominator multiplicities in these degrees are
$$
\smult(\delta_i)=1,\qquad
\smult(\delta_i+\delta_j)=10,\qquad
\smult(2\delta_i+\delta_j)=1,\qquad
\smult(\delta_1+\delta_2+\delta_3)=-64.
$$

\subsection{Constants forced by the Borcherds presentation}

Let \(A=((\delta_i,\delta_j))\), so
$$
A=
\begin{pmatrix}
2&-2&-2\\
-2&2&-2\\
-2&-2&2
\end{pmatrix}.
$$
Choose even Chevalley generators for the real simple roots:
$$
e_i\in\mathfrak g_{\delta_i},\qquad
f_i\in\mathfrak g_{-\delta_i},\qquad
h_i=[e_i,f_i]\qquad(1\leq i\leq3).
$$
The presentation constants are
$$
[h_i,h_j]=0,\qquad
[h_i,e_j]=A_{ij}e_j,\qquad
[h_i,f_j]=-A_{ij}f_j,\qquad
[e_i,f_j]=\delta_{ij}h_i.
$$
Since all \(e_i\) are even,
$$
[e_i,e_i]=0,\qquad [f_i,f_i]=0.
$$
For \(i\neq j\), define the height-two real-commutator vector
$$
E_{ij}:=[e_i,e_j]\in\mathfrak g_{\delta_i+\delta_j}.
$$
Then
$$
E_{ji}=-E_{ij}.
$$
The Cartan constants are
$$
[h_i,E_{ij}]=0,\qquad [h_j,E_{ij}]=0,\qquad
[h_k,E_{ij}]=-4E_{ij}\quad(\{i,j,k\}=\{1,2,3\}).
$$
The lowering constants are forced by Jacobi:
$$
[f_i,E_{ij}]=2e_j,\qquad [f_j,E_{ij}]=-2e_i,\qquad
[f_k,E_{ij}]=0.
$$

For ordered \(i\neq j\), define the height-three real-root vector
$$
E_{iij}:=[e_i,E_{ij}]
=[e_i,[e_i,e_j]]
\in \mathfrak g_{2\delta_i+\delta_j}.
$$
The real Serre relation is exactly
$$
[e_i,E_{iij}]=(\operatorname{ad}e_i)^3e_j=0.
$$
The Cartan constants are
$$
[h_i,E_{iij}]=2E_{iij},\qquad
[h_j,E_{iij}]=-2E_{iij},\qquad
[h_k,E_{iij}]=-6E_{iij}.
$$
The lowering constants are
$$
[f_i,E_{iij}]=2E_{ij},\qquad
[f_j,E_{iij}]=0,\qquad
[f_k,E_{iij}]=0.
$$
All constants in this subsection are denominator-algebra constants.
They are fixed after the displayed Chevalley and commutator
normalizations are chosen.  Rescaling the original simple generators
rescales the named composite vectors, but not the displayed normalized
relations.

\subsection{The first isotropic spaces}

For \(i<j\), set
$$
a_{ij}:=\delta_i+\delta_j.
$$
The real commutator \(E_{ij}\) gives one even vector in
\(\mathfrak g_{a_{ij}}\).  The Gritsenko--Nikulin correction supplies
nine additional even simple imaginary generators on the same primitive
isotropic ray:
$$
u_{ij,1},\ldots,u_{ij,9}\in\mathfrak g_{a_{ij}},
\qquad
\tau(ta_{ij})=9\quad(t\geq1).
$$
Hence
$$
\smult(a_{ij})=1+9=10.
$$
The presentation gives
$$
[h_i,u_{ij,r}]=[h_j,u_{ij,r}]=0,\qquad
[h_k,u_{ij,r}]=-4u_{ij,r}.
$$
Since \(a_{ij}\) is orthogonal to \(\delta_i\) and \(\delta_j\),
the Borcherds orthogonality relations give
$$
[e_i,u_{ij,r}]=[e_j,u_{ij,r}]=0,
\qquad
[f_i,u_{ij,r}]=[f_j,u_{ij,r}]=0.
$$
For the remaining real simple root \(\delta_k\),
$$
(\delta_k,a_{ij})=-4,
$$
so the only real Serre vanishing forced at the simple-generator level is
$$
(\operatorname{ad}e_k)^5u_{ij,r}=0.
$$
The first bracket
$$
[e_k,u_{ij,r}]
\in \mathfrak g_{\delta_1+\delta_2+\delta_3}
$$
is not a numerical constant until a basis in that height-three root
space is chosen.

The isotropic simple generators commute with the real-commutator vector
on the same ray:
$$
[u_{ij,r},E_{ij}]
=[u_{ij,r},[e_i,e_j]]
=[[u_{ij,r},e_i],e_j]+[e_i,[u_{ij,r},e_j]]
=0.
$$
They also commute among themselves on the same ray:
$$
[u_{ij,r},u_{ij,s}]=0.
$$
These vanishings are denominator presentation constants, not Hall
vanishing theorems.

\subsection{The first hyperbolic height-three root}

There are three natural real-simple bracketings in degree
$$
\delta_{123}:=\delta_1+\delta_2+\delta_3=\alpha(1,1,1).
$$
Put
$$
T_1:=[e_1,E_{23}],\qquad
T_2:=[e_2,E_{31}],\qquad
T_3:=[e_3,E_{12}].
$$
Jacobi gives the single forced relation
$$
T_1+T_2+T_3=0.
$$
Moreover
$$
[h_i,T_r]=-2T_r\qquad(1\leq i,r\leq3).
$$
The denominator product gives only the signed total
$$
\smult(\delta_{123})=f(1,1)=-64.
$$
It does not choose a basis of \(\mathfrak g_{\delta_{123}}\), does not
split even and odd dimensions, and does not give the constants expressing
\([e_k,u_{ij,r}]\) or \(T_r\) in such a basis.  This is the first place
where a bracket table cannot be continued from the determinant and the
presentation constants alone.

\subsection{Candidate Hall derivation and failure}

If a compact Hall realization existed, the constants would come from
extension correspondences
$$
\mathfrak M_\gamma^\sigma\times\mathfrak M_\eta^\sigma
\xleftarrow{\ p_{\gamma,\eta}\ }
\mathfrak E_{\gamma,\eta}^\sigma
\xrightarrow{\ q_{\gamma,\eta}\ }
\mathfrak M_{\gamma+\eta}^\sigma
$$
and an oriented pull-push product
$$
m_{\gamma,\eta}
=q_{\gamma,\eta!}p_{\gamma,\eta}^{*}.
$$
For protected primitive bases \(u_{\gamma,a}\), the Hall constants would
be
$$
[u_{\gamma,a},u_{\eta,b}]_{\mathrm{Hall}}
=
\sum_c C^c_{(\gamma,a),(\eta,b)}u_{\gamma+\eta,c}.
$$
The determinant supplies only
$$
\sum_a(-1)^{|a|}
=
\sdim\Prim_{\mathrm{prot},\gamma}
=f(nm,l).
$$
It supplies no moduli stacks \(\mathfrak M_\gamma^\sigma\), no extension
stack \(\mathfrak E_{\gamma,\eta}^\sigma\), no orientation sign, and no
pull-push class.  Therefore it supplies no \(C^c_{(\gamma,a),(\eta,b)}\).

The simplest attempted replacement is the symmetric Euler rule
$$
[x_\gamma,x_\eta]_{\mathrm{trial}}
=\langle\gamma,\eta\rangle_{\mathrm{BPS}}\,x_{\gamma+\eta}.
$$
It fails before geometry enters.  For a real simple charge
\(\gamma_i\),
$$
\langle\gamma_i,\gamma_i\rangle_{\mathrm{BPS}}=-1,
$$
but \(x_{\gamma_i}\) is even, so super-skew symmetry forces
$$
[x_{\gamma_i},x_{\gamma_i}]=0.
$$
For \(i\neq j\), the symmetric rule gives the same scalar in both
orders, while an even Lie bracket requires opposite signs:
$$
[x_{\gamma_i},x_{\gamma_j}]
=-[x_{\gamma_j},x_{\gamma_i}].
$$
Thus the symmetric BPS pairing is the Cartan form after the
anti-isometry \(\alpha\).  It is not the Hall commutator kernel.  Any
honest Hall derivation must retain an antisymmetric Euler form or
orientation cocycle before passing to the three Gram invariants
\((n,l,m)\).

\subsection{Why the determinant cannot recover constants}

\begin{quote}
\textbf{Proposition.}
The Borcherds product for \(64^{-1}\Delta_5(2Z)\), even together with
the active support \(\Gamma_{\mathrm{act}}\), cannot determine the
Hall or denominator structure constants.
\end{quote}

\noindent
\textit{Proof.}
The product records the integers
$$
\smult(\alpha(\gamma))
=\dim\mathfrak g_{\alpha(\gamma),\overline0}
-\dim\mathfrak g_{\alpha(\gamma),\overline1}.
$$
It is unchanged if one replaces a homogeneous root space by another
super vector space with the same signed dimension.  It is also unchanged
if one puts the zero bracket on the same graded super vector spaces.
The zero bracket and the Borcherds bracket have the same determinant but
different structure constants, since in the Borcherds algebra
$$
[e_i,e_j]=E_{ij}\neq0\qquad(i\neq j).
$$
Therefore no operation on the determinant can recover the constants.
\(\square\)

\subsection{Replacement theorem language}

The following theorem is the exact replacement for any manuscript
sentence asserting a constructed microscopic BPS bracket.

\begin{quote}
\textbf{Theorem.}
The Borcherds--Gritsenko--Nikulin construction determines a
charge-graded denominator Lie superalgebra
\[
\mathcal A_{\mathrm{den}}=\mathfrak g_{\Delta_5}
=\bigoplus_{\gamma\in\Gamma_{\mathrm{BPS}}}
\mathcal A_{\mathrm{den},\gamma}
\]
with
\[
[\mathcal A_{\mathrm{den},\gamma},
\mathcal A_{\mathrm{den},\eta}]
\subset
\mathcal A_{\mathrm{den},\gamma+\eta},
\qquad
\operatorname{den}(\mathcal A_{\mathrm{den}})
=64^{-1}\Delta_5(2Z).
\]
After choosing Chevalley generators for the three real simple roots,
the low-degree constants are the Cartan, Chevalley, Serre, and
orthogonality constants displayed in
Section~\ref{sec:bps-bracket-constants-attack-heal}.  These are
denominator-algebra constants.  They are not compact Hall constants.

An identification
\[
\Prim_{\mathrm{prot}}
\mathcal P^{\mathrm{Hall}}_{K3\times E,\sigma,o}
\cong
\mathcal A_{\mathrm{den}}
\]
as Lie superalgebras requires, in addition, a compact oriented critical
Hall correspondence, a protected integration functor, an orientation
cocycle descending to \(\nu_{\Delta_5}\), and a proof that the Hall
primitive commutator satisfies the Borcherds relations above.  The
determinant and scalar square prove only the signed Euler
characteristics; they do not determine the Hall structure constants.
\end{quote}

\subsection{Check table}

\begin{center}
\begingroup
\small
\begin{tabular}{p{0.23\textwidth}|p{0.36\textwidth}|p{0.25\textwidth}}
\textbf{Object} & \textbf{Computed formula} & \textbf{Status}\\
\hline
Cartan matrix &
\(A_{ii}=2,\ A_{ij}=-2\ (i\neq j)\) &
denominator constant\\
Chevalley brackets &
\([h_i,e_j]=A_{ij}e_j,\ [e_i,f_j]=\delta_{ij}h_i\) &
denominator constant\\
height-two real bracket &
\(E_{ij}=[e_i,e_j]=-E_{ji}\) &
constant after naming \(E_{ij}\)\\
height-two lowering &
\([f_i,E_{ij}]=2e_j,\ [f_j,E_{ij}]=-2e_i\) &
Jacobi consequence\\
height-three real root &
\(E_{iij}=[e_i,E_{ij}],\ [e_i,E_{iij}]=0\) &
Serre consequence\\
first isotropic space &
\(\mathfrak g_{\delta_i+\delta_j}\) has signed size \(10=1+9\) &
denominator multiplicity\\
triple real bracketing &
\(T_1+T_2+T_3=0\), \(\smult(\delta_{123})=-64\) &
Jacobi; determinant only\\
trial Euler bracket &
\(\langle\gamma,\eta\rangle_{\mathrm{BPS}}x_{\gamma+\eta}\) &
false: violates super-skew\\
Hall constants &
\(C^c_{(\gamma,a),(\eta,b)}\) &
open: need compact correspondence
\end{tabular}
\endgroup
\end{center}
